\section{Verified exact hulls and finite semidefinite lifts}
\label{sec:formal-exact-hull}

For $r\geq2$, let
\[
 S=\{(u,v)\in\R^r\times\R^r:
 \|u\|^2<1,\ \|v\|^2<1,\ u\transpose v>1/2\},
\]
and let $T$ denote the system obtained by replacing all three inequalities
by weak inequalities. Write $p=1-\|u\|^2$, $q=1-\|v\|^2$, and
$d=u\transpose v$. This Lean package proves
\begin{align*}
 \conv S&=\{(u,v):p>0,\ q>0,\ d+\sqrt{pq}>1/2\},\\
 \cl(\conv S)=\conv T
 &=\{(u,v):p\geq0,\ q\geq0,\ d+\sqrt{pq}\geq1/2\}.
\end{align*}
Every equality includes $r=2$. The proof of sufficiency is a direct
two-point construction, so it does not depend on an unformalized general
good-aggregation hull theorem.

\paragraph{A two-point proof in every required dimension.}
Suppose the strict scalar conditions hold. Put $a=\sqrt p$ and $b=\sqrt q$.
Choose a unit vector $e$ perpendicular to $bu-av$; this requires only
$r\geq2$. Orthogonality implies
\[
 \frac{u\transpose e}{a}=\frac{v\transpose e}{b}=h.
\]
Choose
\[
 \max\left\{0,\frac{1/2-d}{ab}\right\}<k<1,
 \qquad t_\pm=-h\pm\sqrt{h^2+k}.
\]
The two roots have opposite signs and satisfy $2ht_\pm+t_\pm^2=k$.
For $x(t)=(u+tae,v+tbe)$ and $t\in\{t_-,t_+\}$, direct expansion gives
\[
 \|u+tae\|^2=\|u\|^2+pk,\qquad
 \|v+tbe\|^2=\|v\|^2+qk,\qquad
 (u+tae)\transpose(v+tbe)=d+abk.
\]
Thus both $x(t_-)$ and $x(t_+)$ lie in $S$. The positive weights
$t_+/(t_+-t_-)$ and $-t_-/(t_+-t_-)$ sum to one and recover $(u,v)$.
This proves hull sufficiency, including the four-variable case.

Necessity follows from the already verified validity of every good-cone
aggregation on $\conv S$. The coordinate multipliers give $p,q>0$;
the multiplier $(q,p,2\sqrt{pq})$ gives $d+\sqrt{pq}>1/2$.
The package also proves the weak cone-test implication, with explicit
separating multipliers when $p=0$ or $q=0$. It does not infer the weak
claim by replacing strict signs in the strict proof.

\paragraph{Actual affine matrix lifts.}
Define the symmetric matrix
\[
 L(u,v,\sigma)=
 \begin{pmatrix}
 1&\sigma&u\transpose\\
 \sigma&1&v\transpose\\
 u&v&I_r
 \end{pmatrix}.
\]
The formalization proves both the PSD and positive definite
Schur-complement equivalences for this matrix and eliminates $\sigma$.
Consequently,
\begin{align*}
 (u,v)\in\conv S
 &\quad\Longleftrightarrow\quad
 \exists\sigma>1/2:\ L(u,v,\sigma)\succ0,\\
 (u,v)\in\cl(\conv S)
 &\quad\Longleftrightarrow\quad
 \exists\sigma\geq1/2:\ L(u,v,\sigma)\succeq0.
\end{align*}
The latter representation also describes $\conv T$. Singular Schur
complements and both signs of $\sigma-d$ are included.

\paragraph{Closure and the original weak system.}
For a point in the weak scalar region and $0\leq s<1$, the scalar margin
of $(su,sv)$ above $1/2$ is at least $(1-s^2)/2$. This follows from
concavity of the square root of the two-by-two determinant. Thus radial
contraction approximates every weak-region point by strict-hull points.
Continuity gives the reverse closure inclusion.

To prove equality with $\conv T$, the proof uses compactness of $T$ and
of the image of $[0,1]\times T\times T$ under
$(\alpha,y,z)\mapsto\alpha y+(1-\alpha)z$. The two-point construction
places the strict hull inside this compact image, hence its closure lies
there as well. The image is contained in $\conv T$; convexity of the weak
scalar region gives the other inclusion.

Finally, the package identifies $\conv S$ with the intersection of all
strict good aggregations, and $\cl(\conv S)$ with the intersection of all
weak good aggregations. Goodness retains its source definition using the
actual homogeneous eigenvalue count and ordinary-hull validity. These
equalities supply the exact descriptions for which the preceding package
proved that no countable strict good description and no finite weak good
description can suffice.

\paragraph{Scope.}
The frozen claims and declaration coverage
are in \path{supplement/lean/audits/30-infinite-aggregation-hull/}.

These sources extend \path{Formal.InfiniteAggregation} without changing the earlier
certified modules. This package does not verify a numerical SDP solver,
quantitative approximation, arbitrary-quadratic impossibility, or a
general aggregation hull theorem. The accuracy bounds have a separate
verification package.
